\documentclass[11pt,a4paper]{article}
\usepackage[utf8]{inputenc}
\usepackage[T1]{fontenc}
\usepackage[spanish]{babel}
\usepackage{geometry}
\usepackage{listings}
\usepackage{xcolor}
\usepackage{amsmath}
\usepackage{amssymb}
\usepackage{hyperref}
\usepackage{enumitem}

\geometry{margin=2.5cm}

\definecolor{codegray}{rgb}{0.5,0.5,0.5}
\definecolor{codegreen}{rgb}{0,0.6,0}
\definecolor{codeblue}{rgb}{0,0,0.8}
\definecolor{backcolour}{rgb}{0.96,0.96,0.96}

\lstdefinestyle{cstyle}{
    backgroundcolor=\color{backcolour},
    commentstyle=\color{codegreen},
    keywordstyle=\color{codeblue}\bfseries,
    numberstyle=\tiny\color{codegray},
    stringstyle=\color{red},
    basicstyle=\ttfamily\small,
    breakatwhitespace=false,
    breaklines=true,
    captionpos=b,
    keepspaces=true,
    numbers=left,
    numbersep=5pt,
    showspaces=false,
    showstringspaces=false,
    showtabs=false,
    tabsize=2,
    language=C
}

\lstset{style=cstyle}

\title{\textbf{Ejercicios de Pthreads en C}\\ \large De nivel básico a intermedio}
\author{PMD}
\date{\today}

\begin{document}

\maketitle

\section*{Introducción}
Los siguientes ejercicios están diseñados para practicar el uso de hilos POSIX (\texttt{pthreads}) en C. Se cubren tres temas fundamentales:
\begin{itemize}
    \item \textbf{Mutex} (\texttt{pthread\_mutex\_t}): exclusión mutua para proteger secciones críticas.
    \item \textbf{pthread\_join}: sincronización para esperar la terminación de hilos.
    \item \textbf{Variables de condición} (\texttt{pthread\_cond\_t}): comunicación y espera entre hilos.
\end{itemize}

Se recomienda compilar con:
\begin{lstlisting}[language=bash,numbers=none]
gcc archivo.c -o programa -lpthread
\end{lstlisting}

\bigskip

%========================================================
\section{Nivel Básico}
%========================================================

\subsection*{Ejercicio 1: Hola Mundo con Hilos y pthread\_join}
Escribe un programa que cree \textbf{5 hilos}. Cada hilo debe imprimir su identificador (recibido como argumento) y el mensaje \texttt{"Hilo i ejecutándose"}. El hilo principal debe esperar a que \textbf{todos} los hilos terminen usando \texttt{pthread\_join} antes de imprimir \texttt{"Todos los hilos han terminado"}.

\textbf{Objetivo:} Practicar la creación de hilos y \texttt{pthread\_join}.

%--------------------------------------------------------
\subsection*{Ejercicio 2: Suma de un arreglo con mutex}
Dado un arreglo de enteros de tamaño $N = 1000000$ inicializado con valores aleatorios, crea \textbf{4 hilos}. Cada hilo suma una porción del arreglo y acumula el resultado en una variable global \texttt{suma\_total}. Protege el acceso a \texttt{suma\_total} con un \texttt{pthread\_mutex\_t}.

\textbf{Objetivo:} Comprender el uso básico de mutex para evitar condiciones de carrera.

%--------------------------------------------------------
\subsection*{Ejercicio 3: Contador compartido}
Implementa un programa con \textbf{10 hilos}, donde cada hilo incrementa un contador global \textbf{100000 veces}. 
\begin{enumerate}[label=\alph*)]
    \item Primero realiza el incremento \textbf{sin} mutex y observa el resultado final.
    \item Luego repite el experimento \textbf{con} mutex.
\end{enumerate}
Explica la diferencia observada.

\textbf{Objetivo:} Visualizar experimentalmente el problema de las condiciones de carrera.

\bigskip

%========================================================
\section{Nivel Intermedio}
%========================================================

\subsection*{Ejercicio 4: Productor--Consumidor con buffer acotado}
Implementa el clásico problema \textbf{productor--consumidor} usando:
\begin{itemize}
    \item Un buffer circular de tamaño $B = 5$.
    \item \textbf{2 hilos productores} que insertan números enteros aleatorios.
    \item \textbf{2 hilos consumidores} que extraen y muestran los valores.
    \item Un \texttt{pthread\_mutex\_t} para proteger el buffer.
    \item \textbf{Dos variables de condición}: \texttt{no\_lleno} y \texttt{no\_vacio}.
\end{itemize}

Los productores se detienen después de insertar $50$ elementos cada uno. El programa debe terminar cuando todos los elementos hayan sido consumidos.

\textbf{Objetivo:} Combinar mutex y variables de condición para sincronizar hilos.

%--------------------------------------------------------
\subsection*{Ejercicio 5: Barrera de sincronización con variables de condición}
Implementa una \textbf{barrera} manual usando \texttt{pthread\_mutex\_t} y \texttt{pthread\_cond\_t}. La barrera debe tener las siguientes funciones:
\begin{lstlisting}
typedef struct {
    pthread_mutex_t mutex;
    pthread_cond_t cond;
    int count;
    int total;
    int generation;
} barrera_t;

void barrera_init(barrera_t *b, int total);
void barrera_wait(barrera_t *b);
void barrera_destroy(barrera_t *b);
\end{lstlisting}
Crea \textbf{6 hilos} que ejecuten 3 fases de trabajo. Cada hilo debe llamar a \texttt{barrera\_wait()} al final de cada fase, de modo que ningún hilo avance a la siguiente fase hasta que todos hayan terminado la actual.

\textbf{Objetivo:} Entender el uso avanzado de variables de condición y generaciones.

%--------------------------------------------------------
\subsection*{Ejercicio 6: Cola de tareas (thread pool simplificado)}
Implementa un \textbf{pool de 3 hilos trabajadores} que consuman tareas de una cola compartida. Cada tarea es una función simple (por ejemplo, imprimir un número y dormir 100~ms). El programa debe:
\begin{enumerate}[label=\alph*)]
    \item Encapsular la cola con mutex y variables de condición.
    \item Permitir que el hilo principal agregue 20 tareas.
    \item Cuando ya no haya tareas y se indique el cierre (\texttt{shutdown}), los trabajadores deben terminar limpiamente.
\end{enumerate}

\textbf{Objetivo:} Aplicar mutex y variables de condición a un patrón real de concurrencia.


\end{document}
